\documentclass[10.5pt]{article}
\usepackage[a4paper,margin=16mm,top=5mm,bottom=4mm]{geometry}
\usepackage[T1]{fontenc}
\usepackage{helvet}
\renewcommand{\familydefault}{\sfdefault}
\usepackage{xcolor}
\usepackage{hyperref}
\usepackage{enumitem}
\usepackage[brazil]{babel}

\definecolor{acc}{HTML}{2563EB}
\definecolor{textgray}{HTML}{555555}
\definecolor{rule}{HTML}{CBD5E1}

\hypersetup{colorlinks=true,urlcolor=acc,linkcolor=acc,
  pdfauthor={Bernardo Micol Righi},pdftitle={Bernardo Righi Curriculo}}

\setlength\parindent{0pt}\setlength\parskip{0pt}
\pagestyle{empty}

\newcommand\sect[1]{\vspace{3pt}{\color{acc}\bfseries\fontsize{10.5pt}{12pt}\selectfont #1}%
  \par\vspace{0.5pt}{\color{rule}\hrule height .6pt}\vspace{2pt}}

\newcommand\orgline[1]{{\color{acc}\underline{#1}}\par\vspace{0.5pt}}

\newcommand\datesline[1]{{\color{textgray}\itshape\fontsize{8.5pt}{10pt}\selectfont #1}\par\vspace{1.5pt}}

\newenvironment{bullets}{%
  \begin{itemize}[leftmargin=4mm,label={\color{acc}\fontsize{8pt}{8pt}\selectfont --},itemsep=0.5pt,topsep=0pt,parsep=0pt]
}{\end{itemize}\vspace{0pt}}

\newcommand\projtitle[3]{% title, stack, github-url-text
  {\bfseries #1}\par
  {\fontsize{9pt}{11pt}\selectfont #2 \textbullet\ \href{https://#3}{#3}}\par\vspace{0.5pt}}

\begin{document}

{\fontsize{19pt}{21pt}\selectfont\bfseries Bernardo Micol Righi}\par\vspace{1pt}
{\color{acc}\itshape\fontsize{10pt}{12pt}\selectfont Desenvolvedor Back-End \textperiodcentered\ Sistemas de IA \textperiodcentered\ Tech Influencer}\par\vspace{3pt}
{\fontsize{9pt}{11pt}\selectfont
+55 (51) 99601-1501 \textperiodcentered\ \href{mailto:bernardomicolrighi@outlook.com}{bernardomicolrighi@outlook.com} \textperiodcentered\ \href{https://righi.dev}{righi.dev} \textperiodcentered\ \href{https://github.com/righibe}{github.com/righibe}\par
\href{https://linkedin.com/in/bernardo-righi}{linkedin.com/in/bernardo-righi} \textperiodcentered\ \href{https://lattes.cnpq.br/8608406696939587}{lattes.cnpq.br/8608406696939587}
}

\sect{SOBRE}
{\fontsize{9.3pt}{12pt}\selectfont
Desenvolvedor back-end com 2 anos de experiencia em Python, IA e automacao. Pesquisador na Unisinos em parceria
com a Dell, focado em observabilidade de agentes LLM usando LangChain e OpenTelemetry. Co-fundador do \href{https://discord.gg/}{Servidor dos Programadores}, uma das
top-6 comunidades dev do Brasil no Discord (18 mil+ membros). Tambem crio conteudo sobre programacao e IA como
tech influencer. Fluente em ingles.}

\sect{EXPERIENCIA PROFISSIONAL}

{\bfseries\fontsize{10pt}{12pt}\selectfont Pesquisador --- Agentes de IA \& Observabilidade}\par\vspace{1pt}
\orgline{Unisinos / Dell Technologies}
\datesline{Jun 2026 -- Atual \textperiodcentered\ Remoto}
\begin{bullets}
  \item\fontsize{9.2pt}{11.7pt}\selectfont Pesquisando padroes de observabilidade para sistemas multiagente, estudando como rastrear comportamento, latencia e falhas de agentes de IA em producao.
  \item\fontsize{9.2pt}{11.7pt}\selectfont Investigando estrategias para rastrear chamadas de ferramentas, transferencias entre agentes (handoffs) e propagacao de contexto --- projeto em fase de pesquisa e prototipagem de arquitetura junto ao time da Dell.
  \item\fontsize{9.2pt}{11.7pt}\selectfont Entreguei um estudo comparativo de 6 pipelines de clusterizacao por embeddings (K-Means, UMAP, PaCMAP, LocalMAP, densMAP + HDBSCAN) sobre dados reais de atendimento, com validacao estatistica (trustworthiness, DBCV) fundamentada na literatura cientifica, orientando a abordagem do time para descoberta automatica de topicos.
\end{bullets}

{\bfseries\fontsize{10pt}{12pt}\selectfont Co-fundador \& Desenvolvedor Back-End}\par\vspace{1pt}
\orgline{Servidor dos Programadores}
\datesline{Jun 2022 -- Atual \textperiodcentered\ Remoto}
\begin{bullets}
  \item\fontsize{9.2pt}{11.7pt}\selectfont A comunidade crescia rapido, mas perdia membros por onboarding e moderacao fracos --- construi bots em Python/Discord.py que automatizaram boas-vindas e suporte tecnico, melhorando a retencao e reduzindo o trabalho manual do time.
  \item\fontsize{9.2pt}{11.7pt}\selectfont Com o crescimento acelerado, a infraestrutura passou a exigir deploys sem downtime --- montei pipelines de CI/CD com GitHub Actions que testam e publicam atualizacoes automaticamente a cada push, eliminando deploys manuais.
  \item\fontsize{9.2pt}{11.7pt}\selectfont Liderei a organizacao de eventos tecnicos com patrocinadores como Alura, Shard Cloud e Hostinger, coordenando palestrantes e sessoes ao vivo para 18 mil+ desenvolvedores.
\end{bullets}

\sect{PROJETOS}

\projtitle{Pipeline de CI/CD em Python}{Python \textperiodcentered\ GitHub Actions \textperiodcentered\ Docker \textperiodcentered\ Flask}{github.com/righibe/ci-cd-pipeline-python}
\begin{bullets}
  \item\fontsize{9.2pt}{11.7pt}\selectfont Situacao: nao havia processo automatizado de testes ou deploy --- cada atualizacao era um risco de bug em producao.
  \item\fontsize{9.2pt}{11.7pt}\selectfont Acao: construi um pipeline completo com GitHub Actions e Docker, rodando testes automaticos a cada push e fazendo deploy no container so se tudo passasse.
  \item\fontsize{9.2pt}{11.7pt}\selectfont Resultado: eliminei os deploys manuais e ganhei a garantia de que so codigo validado chegava ao ambiente.
\end{bullets}

\projtitle{Predicao de Precos de Imoveis --- Sao Paulo}{Python \textperiodcentered\ scikit-learn \textperiodcentered\ Decision Tree \textperiodcentered\ Random Forest}{github.com/righibe/predict-prices-sp}
\begin{bullets}
  \item\fontsize{9.2pt}{11.7pt}\selectfont Situacao: queria aplicar ML supervisionado a um problema real, indo alem da teoria.
  \item\fontsize{9.2pt}{11.7pt}\selectfont Acao: treinei e comparei Decision Tree e Random Forest com scikit-learn para prever precos de imoveis em Sao Paulo, analisando a importancia das features.
  \item\fontsize{9.2pt}{11.7pt}\selectfont Resultado: identifiquei localizacao e metragem como as variaveis dominantes e consolidei minha base pratica em regressao.
\end{bullets}

\projtitle{API REST de Cadastro de Pessoas --- Java}{Java \textperiodcentered\ Spring Boot \textperiodcentered\ JPA / Hibernate \textperiodcentered\ Spring Web}{github.com/righibe/Cadastro-pessoas}
\begin{bullets}
  \item\fontsize{9.2pt}{11.7pt}\selectfont Situacao: eu so tinha experiencia em Python e precisava entender o ecossistema Java.
  \item\fontsize{9.2pt}{11.7pt}\selectfont Acao: construi uma API REST do zero com Spring Boot e JPA, seguindo o fluxo Controller $\rightarrow$ Service $\rightarrow$ Repository com DTO e ORM.
  \item\fontsize{9.2pt}{11.7pt}\selectfont Resultado: uma base solida em Java e clareza sobre o padrao MVC para projetos corporativos.
\end{bullets}

\sect{HABILIDADES TECNICAS}
{\fontsize{9.2pt}{11.7pt}\selectfont
\begin{tabular}{@{}p{32mm}p{140mm}@{}}
\textbf{Linguagens} & Python, Java, C, C++, JavaScript, SQL \\
\textbf{IA \& Agentes} & LangChain, LangGraph, LCEL, RAG, LLMs, scikit-learn \\
\textbf{Observabilidade} & OpenTelemetry, Jaeger, Grafana, agent tracing \\
\textbf{DevOps \& Infra} & Docker, GitHub Actions, CI/CD, Git, Linux \\
\textbf{Frameworks} & Spring Boot, Django, Flask, Discord.py \\
\textbf{Bancos de Dados} & MySQL, JPA / Hibernate \\
\end{tabular}}

\sect{FORMACAO}
{\bfseries\fontsize{10pt}{12pt}\selectfont Bacharelado em Ciencia da Computacao}\par\vspace{1pt}
{\fontsize{8.7pt}{11pt}\selectfont\itshape Universidade do Vale do Rio dos Sinos (Unisinos)\/\ \textperiodcentered\ Fev 2026 -- Atual \textperiodcentered\ Sao Leopoldo, Brasil \textperiodcentered\ \href{https://lattes.cnpq.br/8608406696939587}{lattes.cnpq.br/8608406696939587}}

\sect{CURSOS \& CERTIFICACOES}
{\fontsize{9.2pt}{11.7pt}\selectfont
Maratona Java (Java Virado no Jiraya) \textperiodcentered\ Python Mundo 1, 2 e 3 \textperiodcentered\ HTML5 \& CSS3 (Curso em Video) \textperiodcentered\ Imersao DevOps \&
Python (Alura) \textperiodcentered\ Design and Social Enterprise HDSE Leadership (Lehigh University)}

\sect{IDIOMAS}
{\fontsize{9.2pt}{11.7pt}\selectfont
\textbf{Portugues} --- Nativo ~~$\vert$~~ \textbf{Ingles} --- Fluente (C1) ~~$\vert$~~ \textbf{Espanhol} --- Intermediario (B1)}

\vfill
{\color{rule}\hrule height .5pt}\vspace{2.5mm}
{\color{textgray}\fontsize{8pt}{10pt}\selectfont
\href{https://righi.dev}{righi.dev} \textperiodcentered\ \href{https://github.com/righibe}{github.com/righibe} \textperiodcentered\ \href{https://linkedin.com/in/bernardo-righi}{linkedin.com/in/bernardo-righi} \textperiodcentered\ \href{https://lattes.cnpq.br/8608406696939587}{lattes.cnpq.br/8608406696939587} \textperiodcentered\ Atualizado em junho de 2026}

\end{document}
